For an instance of mix with mix formula \(A\) and immediate
subderivations \(D_1\) and \(D_2\), its rank is \(r(A)\) and its height
is \(h(D_1)+h(D_2)\), the sum of the heights of those subderivations.
The reduction is controlled by the lexicographic order on
\[
  \tuple{r(A),h(D_1)+h(D_2)}.
\]
For a cut, the mix instance \(m=n=1\), these measures are its \emph{cut
rank} and \emph{cut height}.

\begin{theorem}[Cut elimination for \(\LK\) and \(\LJ\)]
Every \(\LK\)-derivation can be transformed into a cut-free
\(\LK\)-derivation of the same endsequent.  Every \(\LJ\)-derivation can be
transformed into a cut-free \(\LJ\)-derivation of the same endsequent.
\textit{Status: proof sketch.}
\end{theorem}

\paragraph{Proof sketch.}
By \cref{lem:cut-mix-equipollence}, first replace every cut in the given
derivation by the mix instance \(m=n=1\); the reduction below eliminates
mixes, and steps 1, 2, 5, and 6 apply verbatim with ``mix'' for ``cut''
and the mix formula for the cut formula.
Choose an uppermost cut of maximal cut rank.
\begin{enumerate}
  \item If the cut formula is not principal in the final inference of one
        premise, and that inference is not a contraction of the cut formula,
        commute the cut above that inference.  The cut rank is
        unchanged and the cut height decreases.
  \item If the cut formula is introduced by weakening in one premise, delete
        the cut and restore the conclusion by weakenings and contractions.
  \item If a premise of a mix ends in a contraction of the mix formula
        (\textsc{CL} above the right premise, \textsc{CR} above the
        left), replace that premise by the contraction's premise and
        raise the corresponding multiplicity by one: the rank is
        unchanged and the height decreases by one.
  \item If the cut formula is principal on both sides, replace the cut by cuts
        on proper subformulas.  For example, a cut between \(\land R\) and
        \(\land L\) is replaced by cuts on the two conjuncts.  Every new cut
        has strictly smaller cut rank, so the number of cuts of maximal cut
        rank decreases.  When the mix formula is principal on both sides
        with multiplicities \((m,n)\), the principal reduction yields
        mixes on proper subformulas, of smaller rank, together with
        residual mixes on the remaining occurrences, of equal rank but
        smaller height.
  \item If an initial sequent is a premise of the cut, delete the cut directly.
  \item For quantifiers, rename eigenvariables before permutation and use the
        substitution lemma.  This prevents capture and preserves the
        eigenvariable conditions.
  \item Repeat according to the well-founded lexicographic measure.  The
        resulting derivation has no cuts.
\end{enumerate}
The calculus-specific table of commuting and principal reductions is the
omitted verification; it occupies the
complete proof in
\cite{TroelstraSchwichtenberg2000,NegriVonPlato2001}.

